\section{Introduction}
\label{sec:intro}

Software bugs are unavoidable, but how much they cost depends on how quickly and
accurately a team can classify them. Defect classification is what turns a raw
bug report into that identification: different defect types call for
different debugging strategies, resource allocations, and repair approaches
\citep{catolino2019not,hirsch2022using}. Manual triage still gives the best results, but it does not scale to large
benchmark datasets \citep{andrade2025empirical}, so we need a dataset large
enough to test automated classification on, and a benchmark rich enough to
classify from.

\textit{Defects4J} \citep{just2014defects4j} is exactly that benchmark, the one
most empirical Java fault studies already build on. It gives us a large,
curated set of reproducible bugs across 17 open-source projects, and every bug
comes with a rich bundle of evidence: the original bug report, the tests that
trigger it, the buggy and fixed source, and the buggy-to-fixed diff. That
richness is exactly what makes \textit{Defects4J} such a good fit for automated
classification research.

And yet, despite this richness, no one has produced a full, semantic
labelling of the benchmark. We looked at why, and existing classification work
falls into three groups, each with a different focus: control- and data-flow
analyses capture \emph{how} a bug manifests syntactically
\citep{vanderspuy2025anatomy}; patch taxonomies describe \emph{what} developers
changed \citep{sobreira2018dissection}; and small-scale manual inspections do
not scale \citep{xuan2017better}. None of them classifies bugs into a semantic
defect-type taxonomy at the scale of the full benchmark.

So where does that leave us? LLMs have shown that code-aware reasoning can
classify software artefacts competitively, as long as the prompt is designed
well \citep{colavito2024llm,koyuncu2025exploring,nong2024cot}. Left
unconstrained, though, LLM classification suffers from label drift and
hallucination \citep{gao2024rcegen}: with no fixed label space, different runs
assign incompatible categories to the same bug. This can be solved with
Orthogonal Defect Classification (ODC), a taxonomy that classifies each defect
into one of \textbf{seven} Design/Code defect types capturing the semantics of the bug itself \citep{chillarege1992odc,ibm2013odc}.
Pairing LLMs with ODC fixes that by turning classification into
a bounded, semantically grounded decision problem, but a forced, fixed schema
brings back the opposite failure: it manufactures false certainty for defects
that do not fit. Our answer is to treat the taxonomy's own coverage as
something to measure rather than assume, by giving the model an explicit escape
category \textbf{"Other":} a defect that falls outside the seven types becomes recorded data
instead of a mislabelled forced choice.

\subsection{Research Questions}
\label{sec:rqs}

Turning that idea into something we can actually test means asking five
concrete questions. We organise this work around them:

\begin{description}
  \item[RQ1: Defect-type distribution across projects.] What is the
        distribution of ODC defect types across the \textit{Defects4J} benchmark,
        and does this distribution vary significantly across projects?
  \item[RQ2: Coverage of the ODC taxonomy.] How much of \textit{Defects4J} is
        covered by the seven standard ODC defect types, and how often does a
        bug instead need the \textbf{Other} escape category?
  \item[RQ3: Classification accuracy.] How well does the LLM-based pipeline,
        using the enforced scientific reasoning loop and the ODC taxonomy,
        classify \textit{Defects4J} bugs into the correct defect types under a
        four-level accuracy evaluation?
  \item[RQ4: Contribution of each component.] How much do the scientific
        reasoning loop and explicit ODC taxonomy grounding each improve
        classification accuracy, label consistency, and vocabulary reduction
        relative to an unstructured LLM baseline?
  \item[RQ5: Effect of seeing the fix.] What is the magnitude and pattern of
        semantic divergence between pre-fix and post-fix ODC classifications, and
        how does per-project classification reliability vary across projects?
\end{description}

RQ1 and RQ2 characterise the dataset and the taxonomy's fit to it; RQ3 measures
how accurately our default configuration classifies; RQ4 isolates what the
structured components contribute; and RQ5 quantifies how much the available
evidence, not just the bug, shapes the assigned type.

\subsection{Contributions}
\label{sec:contributions}

Our main contributions are:
\begin{itemize}
  \item We separate \emph{what label space} the model uses (the taxonomy) from
        \emph{how it reasons} toward a label (the strategy), giving a small set
        of controlled, comparable conditions instead of one all-in-one prompt.
  \item We give ODC an open ``Other'' escape category that requires a written
        justification, which turns the taxonomy's coverage into something we
        measure directly instead of assume.
  \item We build an enforced scientific reasoning loop, adapted from fault
        localization to defect-type classification, where the model has to
        commit to a hypothesis and probe held-back evidence before it is
        allowed to conclude, leaving behind a full, auditable reasoning
        transcript.
  \item We classify every bug in two evidence modes, pre-fix and post-fix,
        mirroring ODC's own opener/closer split. This lets us check the
        pipeline's self-consistency without ground-truth labels, and lets us
        use the fix-independent Impact attribute (Section~\ref{sec:impact}) as
        a check on how much of the pre-fix/post-fix drift is genuine versus
        just noise.
\end{itemize}

%% =============================================================================
